\chapter*{Resumen}
\addcontentsline{toc}{chapter}{Resumen}

El presente trabajo desarrolla un sistema de gestión de rutas óptimas para reparto de pedidos en la ciudad del Cusco, implementado como aplicación de escritorio en Python con interfaz gráfica basada en Tkinter. El problema central consiste en determinar, de forma eficiente, qué pedidos debe atender un repartidor y en qué orden visitarlos, minimizando la distancia total recorrida y respetando restricciones de capacidad del vehículo y calles cerradas en la red vial. Para resolverlo se implementaron, desde cero y sin librerías de optimización externas, cuatro paradigmas algorítmicos: (1) un algoritmo \textit{Greedy} de vecino más cercano para construcción de ruta; (2) \textit{Merge Sort} (Divide y Vencerás) para ordenamiento de pedidos; (3) \textit{Knapsack 0/1} y \textit{Held-Karp} (Programación Dinámica) para selección óptima de pedidos y ruta TSP exacta; y (4) \textit{Backtracking} con poda de ciclos para exploración de rutas alternativas con calles bloqueadas. Los \textit{benchmarks} empíricos confirmaron la complejidad teórica: Held-Karp alcanzó 42.2 ms para $n=12$ pedidos con un speedup del Greedy de $138.8\times$, mientras que el gap de suboptimalidad del Greedy creció de 0\% en $n=3$ a 14.81\% en $n=12$. El Backtracking encontró entre 3 y 9 rutas alternativas en menos de 1.1 ms en cuatro escenarios del grafo de Cusco. El sistema resultante es estable, visualmente verificable y permite reproducir todos los escenarios de prueba de forma interactiva.

\textit{Palabras clave:} rutas óptimas, algoritmos greedy, programación dinámica, backtracking, Held-Karp, Cusco, TSP, Knapsack
